\documentclass[12pt,oneside]{book}

% ============================================================
% METADATOS
% ============================================================
\newcommand{\universidad}{Universidad de Buenos Aires (UBA)}
\newcommand{\facultad}{Facultad de Derecho}
\newcommand{\materia}{Práctica Profesional}
\newcommand{\titulo}{Apuntes de Práctica Profesional}
\newcommand{\autor}{Isaac Edgar Camacho Ocampo}
\newcommand{\anio}{2026}

% ============================================================
% PAQUETES
% ============================================================
\usepackage[utf8]{inputenc}
\usepackage[T1]{fontenc}
\IfFileExists{spanish.ldf}{%
    \usepackage[spanish,es-nodecimaldot]{babel}%
}{}
\usepackage[
    top=2.5cm,
    bottom=2.5cm,
    left=3cm,
    right=2.5cm,
    headheight=15pt
]{geometry}
\usepackage{amsmath}
\usepackage{amssymb}
\usepackage{mathtools}
\usepackage{graphicx}
\usepackage{float}
\usepackage{caption}
\usepackage{subcaption}
\usepackage{booktabs}
\usepackage{array}
\usepackage{longtable}
\usepackage{tabularx}
\usepackage{enumitem}
\usepackage{xcolor}
\usepackage[most]{tcolorbox}
\usepackage{fancyhdr}
\usepackage{ragged2e}
\usepackage[
    colorlinks=true,
    linkcolor=blue!50!black,
    citecolor=green!40!black,
    urlcolor=blue!60!black,
    pdftitle={\titulo},
    pdfauthor={\autor},
    pdfsubject={Apuntes universitarios},
    bookmarks=true
]{hyperref}

% ============================================================
% CONFIGURACIÓN
% ============================================================
\graphicspath{
    {./img/}
    {./graficos/}
    {./figuras/}
}
\setcounter{tocdepth}{2}
\setlength{\parindent}{0pt}
\setlength{\parskip}{0.5em}
\setlist[itemize]{
    topsep=0.3em,
    itemsep=0.2em
}
\setlist[enumerate]{
    topsep=0.3em,
    itemsep=0.2em
}
\clubpenalty=10000
\widowpenalty=10000

% ============================================================
% COMANDOS
% ============================================================
\newcommand{\abs}[1]{\left|#1\right|}
\newcommand{\norm}[1]{\left\lVert#1\right\rVert}
\newcommand{\vect}[1]{\mathbf{#1}}
\newcommand{\deriv}[2]{\frac{d#1}{d#2}}
\newcommand{\pderiv}[2]{\frac{\partial#1}{\partial#2}}

% ============================================================
% ENTORNOS / CAJAS
% ============================================================
\newtcolorbox{definition}[1]{
    enhanced, breakable,
    title={Definición: #1},
    fonttitle=\bfseries,
    colback=blue!3, colframe=blue!50!black,
    boxrule=0.8pt, arc=2mm
}
\newtcolorbox{important}{
    enhanced, breakable,
    title={Importante},
    fonttitle=\bfseries,
    colback=yellow!8, colframe=orange!70!black,
    boxrule=0.8pt, arc=2mm
}
\newtcolorbox{example}[1]{
    enhanced, breakable,
    title={Ejemplo: #1},
    fonttitle=\bfseries,
    colback=green!4, colframe=green!40!black,
    boxrule=0.8pt, arc=2mm
}
\newtcolorbox{formula}[1]{
    enhanced, breakable,
    title={Fórmula / Regla: #1},
    fonttitle=\bfseries,
    colback=purple!4, colframe=purple!50!black,
    boxrule=0.8pt, arc=2mm
}
\newtcolorbox{exercise}[1]{
    enhanced, breakable,
    title={Ejercicio: #1},
    fonttitle=\bfseries,
    colback=gray!5, colframe=gray!60!black,
    boxrule=0.8pt, arc=2mm
}
\newtcolorbox{note}{
    enhanced, breakable,
    title={Nota},
    fonttitle=\bfseries,
    colback=white, colframe=gray!50,
    boxrule=0.6pt, arc=2mm
}

% ============================================================
% CONFIGURACIÓN FINAL (ENCABEZADOS Y PIES)
% ============================================================
\pagestyle{fancy}
\fancyhf{}
\fancyhead[L]{\nouppercase{\leftmark}}
\fancyhead[R]{\materia}
\fancyfoot[C]{\thepage}
\renewcommand{\headrulewidth}{0.4pt}
\renewcommand{\footrulewidth}{0pt}
\fancypagestyle{plain}{
    \fancyhf{}
    \fancyfoot[C]{\thepage}
    \renewcommand{\headrulewidth}{0pt}
}

% ============================================================
\begin{document}

% ---------------- PORTADA ----------------
\begin{titlepage}
\thispagestyle{empty}
\begin{center}
\vspace*{1cm}
{\large \textsc{\universidad}}

\vspace{0.2cm}
{\large \textsc{\facultad}}

\vspace{3cm}
{\Huge \textbf{\materia}}

\vspace{1cm}
{\LARGE \titulo}

\vspace{0.8cm}
\textit{Comprender los procedimientos es el primer paso para saber actuar cuando el estudio se vuelve práctica.}

\vspace{1.2cm}
\rule{0.70\textwidth}{0.5pt}

\vspace{0.5cm}
{\large Apuntes de estudio}

\vspace{0.5cm}
\rule{0.70\textwidth}{0.5pt}

\vfill
{\large \autor}

\vspace{0.3cm}
{\large \anio}
\end{center}

\vfill
\begin{flushleft}
\textit{Este es un modesto aporte a los estudiantes de ``\materia'' no pretende sustituir el material oficial de la materia.}
\end{flushleft}
\end{titlepage}

% ---------------- ÍNDICE ----------------
\tableofcontents
\clearpage

% ---------------- CONTENIDO ----------------
\chapter[Violencia doméstica]{Protocolo de actuación ante la violencia doméstica}

\section{El planteo: una consulta por violencia doméstica}

Un problema central de la práctica profesional es saber qué hacer cuando una persona consulta por violencia doméstica. La consulta puede presentarse en cualquier ámbito: una clienta o consultante, un familiar, una amiga o cualquier persona cercana que manifieste que en su casa la golpean y que ``no aguanta más''. La persona puede acercarse a un abogado o abogada en la calle, en el trabajo o en el entorno familiar, planteando una situación muy violenta, el deseo de irse o de que el agresor se vaya, y preguntando cómo solucionarlo.

\begin{example}{La consulta inesperada}
Una persona consulta a un abogado o abogada: ``Mirá, estoy en una situación muy violenta, me quiero ir, quiero que se vaya. ¿Cómo lo soluciono?''. Quien recibe la consulta debe saber qué hacer, incluso si la persona ya tenía una medida previa o una denuncia penal.
\end{example}

Conocer el procedimiento resulta relevante porque se trata de una situación que se atiende con frecuencia y que suele presentarse fuera del ámbito profesional; por ello es conveniente saber qué se hace también para orientar a conocidos y amistades. Los procedimientos disponibles contemplan, además, medidas como el botón antipánico.

\section{El recorrido de una denuncia en la Ciudad}

\begin{definition}{Oficina de Violencia Doméstica (OVD)}
Oficina que funciona las 24 horas (Lavalle 1250), donde un \textbf{equipo interdisciplinario} atiende a las personas, toma la denuncia y evalúa el \textbf{nivel de riesgo} de la situación.
\end{definition}

El procedimiento se desarrolla del siguiente modo. La persona es derivada a la Oficina de Violencia Doméstica. Allí, el equipo interdisciplinario toma la denuncia y la vincula con el grado de riesgo, que se evalúa previamente. El legajo se remite directamente al \textbf{juzgado civil}. Cuando el nivel de riesgo es \textbf{medio, medio-alto o superior}, el juzgado civil dicta una \textbf{orden de restricción}. Esa orden se envía a la \textbf{policía}, que se presenta en el domicilio y retira de la casa a la persona violenta; si resulta necesario, deja además una consigna.

\begin{important}
La orden de restricción se dicta cuando el nivel de riesgo es medio, medio-alto o superior. El legajo va directamente de la OVD al juzgado civil, y la ejecución de la orden queda a cargo de la policía.
\end{important}

\subsection{Cuadro: el recorrido de una denuncia}

\begin{table}[H]
\centering
\small
\begin{tabularx}{\textwidth}{
    >{\raggedright\arraybackslash}p{2.6cm}
    >{\raggedright\arraybackslash}X
    >{\raggedright\arraybackslash}p{3.3cm}
}
\toprule
\textbf{Paso} & \textbf{Qué ocurre} & \textbf{Quién interviene} \\
\midrule
\textbf{1. Consulta} & La persona consulta en la calle, en el trabajo o en la familia: ``me quiero ir, quiero que se vaya''. & Abogado/a, familiar o conocido \\
\textbf{2. Derivación} & Se la envía a la Oficina de Violencia Doméstica, abierta las 24 horas; también si ya había una medida o una denuncia penal. & Quien recibe la consulta \\
\textbf{3. Evaluación} & Un equipo interdisciplinario evalúa el nivel de riesgo de la situación. & OVD \\
\textbf{4. Remisión} & El legajo va directamente al juzgado civil. & OVD $\rightarrow$ juzgado civil \\
\textbf{5. Decisión} & Si el riesgo es medio, medio-alto o superior, se dicta una orden de restricción. & Juzgado civil \\
\textbf{6. Ejecución} & La orden llega a la policía, que se presenta en el domicilio y saca a la persona violenta de la casa. & Policía \\
\textbf{7. Medidas} & Según el caso: consigna, botón de pánico y alimentos. & Juzgado / policía \\
\bottomrule
\end{tabularx}
\end{table}

\section{Medidas que pueden acompañar la restricción}

Según la situación, se adoptan distintas medidas. Además de la restricción y del retiro de la persona violenta, la policía puede dejar una \textbf{consigna} en el domicilio. También es probable que se disponga un \textbf{botón de pánico} y que se fijen \textbf{alimentos}. % TODO: contenido incompleto o ambiguo en las notas originales (el apunte menciona ``alimentos'' y ``provisorios'' sin precisar su relación)
En definitiva, las medidas dependen de cada caso.

\section{Denunciar fuera de la Ciudad}

Si la persona afectada o su familiar o conocido vive en la provincia (por ejemplo, en Catán), puede acudir a \textbf{cualquier fiscalía o a la policía más cercana} y realizar la denuncia de la misma manera. La policía debe intervenir en muchos casos y puede llamar a alguien de la OVD para que tome la denuncia, dándole intervención. % TODO: contenido incompleto o ambiguo en las notas originales (el docente expresa ``creo'' respecto de que la policía pueda llamar a alguien de la OVD)
Algunas provincias o municipios tienen convenio con la OVD y cuentan con oficinas propias (por ejemplo, Moreno); en esos casos el procedimiento es ``casi lo mismo''. Si no es así, se acude a la fiscalía de la provincia.

\begin{table}[H]
\centering
\small
\begin{tabularx}{\textwidth}{
    >{\raggedright\arraybackslash}p{3cm}
    >{\raggedright\arraybackslash}X
    >{\raggedright\arraybackslash}X
}
\toprule
\textbf{Criterio} & \textbf{En la Ciudad (Capital)} & \textbf{En la provincia} \\
\midrule
\textbf{Dónde se denuncia} & Oficina de Violencia Doméstica, abierta las 24 horas. & Cualquier fiscalía o la policía más cercana. \\
\textbf{Quién atiende} & Un equipo interdisciplinario que toma la denuncia. & La fiscalía o la policía; esta puede llamar a alguien de la OVD o darle intervención. \\
\textbf{Particularidad} & El legajo va directamente al juzgado civil. & Algunos municipios tienen convenio u oficinas de la OVD (por ejemplo, Moreno): ``casi lo mismo''. \\
\bottomrule
\end{tabularx}
\end{table}

\section{Tipos de violencia}

La violencia no es solamente física, es decir, no se limita a los golpes. Existen muchos tipos de violencia, entre ellos la \textbf{violencia económica} y la \textbf{violencia psicológica}.

\begin{important}
No hay que aguantarse ningún tipo de violencia, porque \textbf{siempre puede escalar}.
\end{important}

\section{Violencia en las causas y lectura de expedientes}

Todos los grupos de trabajo tienen causas de violencia entre las que llevan. Puede suceder, además, que una causa comience como un provisorio y luego se vincule con violencia, o a la inversa. % TODO: contenido incompleto o ambiguo en las notas originales (pasaje poco claro sobre ``provisorios que vienen de violencia'')

Se recomienda leer completos los expedientes de violencia, aun cuando sean largos, y comprender cómo funcionan las cosas. Conocer el funcionamiento del trámite permite \textbf{destrabar} situaciones: cuando se responde que el expediente ``está en despacho'' y que ``después se hace'', quien sabe lo que corresponde hacer puede exigirlo y resolver el problema, ya que los abogados son quienes solucionan problemas.

\begin{note}
Complementariamente, se indicó que el conjunto de las 32 causas que lleva el curso será objeto de comprobación y que todos deben conocerlas, no solo la propia. También se recomendó pasar los documentos a PDF en lugar de imprimirlos como imágenes, porque así se leen mejor y se ahorra tinta.
\end{note}

% ---------------- CORNELL ----------------
\chapter[Repaso Cornell]{Repaso con método Cornell}

\small
\begin{longtable}{
    >{\raggedright\arraybackslash}p{\dimexpr0.28\textwidth-2\tabcolsep\relax}
    >{\raggedright\arraybackslash}p{\dimexpr0.72\textwidth-2\tabcolsep\relax}
}
\toprule
\textbf{Preguntas / palabras clave} &
\textbf{Notas / respuestas} \\
\midrule
\endhead

\textbf{¿Qué debe saber hacer quien recibe una consulta por violencia doméstica?} &
Debe saber derivar a la persona a la Oficina de Violencia Doméstica, abierta las 24 horas, y conocer el recorrido completo del trámite, incluso si ya existía una medida o una denuncia penal. \\[0.4em]

\textbf{¿Qué es la OVD y cómo atiende?} &
Funciona las 24 horas (Lavalle 1250). Un equipo interdisciplinario toma la denuncia y evalúa el nivel de riesgo de la situación. \\[0.4em]

\textbf{¿Cuál es el recorrido posterior a la evaluación de riesgo?} &
El legajo va directamente al juzgado civil. Si el riesgo es medio, medio-alto o superior, el juzgado dicta una orden de restricción, que se envía a la policía; esta se presenta en el domicilio y retira a la persona violenta. \\[0.4em]

\textbf{¿Qué medidas pueden acompañar la restricción?} &
Consigna policial si es necesaria, botón de pánico y alimentos. Las medidas dependen de la situación. \\[0.4em]

\textbf{¿Dónde se denuncia fuera de la Ciudad?} &
En cualquier fiscalía o en la policía más cercana. La policía puede llamar a la OVD o darle intervención; algunos municipios, como Moreno, tienen convenio u oficinas de la OVD. \\[0.4em]

\textbf{¿La violencia es solo física?} &
No. Existen también la violencia económica y la psicológica. No hay que aguantar ninguna, porque siempre puede escalar. \\[0.4em]

\textbf{¿Por qué conviene leer los expedientes de violencia completos?} &
Porque entender cómo funcionan permite destrabar trámites y no conformarse con respuestas como ``está en despacho, después lo hacemos''. Todos los grupos tienen causas de violencia. \\[0.4em]

\textbf{¿Qué se comprobará respecto de las causas del curso?} &
Que todos conozcan las 32 causas que lleva el curso, no solo la propia. \\

\midrule

\multicolumn{2}{p{\dimexpr\textwidth-2\tabcolsep\relax}}{
\textbf{Resumen:}
Ante una consulta por violencia doméstica se deriva a la Oficina de Violencia Doméstica (24 horas), donde se evalúa el riesgo; el legajo pasa al juzgado civil, que dicta una orden de restricción si el riesgo es medio, medio-alto o superior, y la policía la ejecuta, con medidas como consigna, botón de pánico y alimentos. Fuera de la Ciudad se denuncia en fiscalía o policía. La violencia no es solo física y siempre puede escalar; leer los expedientes completos permite destrabar trámites.
} \\

\bottomrule
\end{longtable}

\end{document}
